\begin{abstract}
As urban transit networks digitize, the accurate detection of impending service instabilities—such as bus bunching—requires bridging the gap between raw telemetry and actionable predictive intelligence. Existing Early Warning Systems (EWS) frequently struggle with the inherent discontinuity of General Transit Feed Specification Real-Time (GTFS-RT) data streams, where API rate limits and packet loss degrade feature fidelity. This paper proposes a spatio-temporal sequential framework optimized for Transit Digital Twin architectures that fundamentally reconceptualizes the handling of polling uncertainty. The primary contribution is a novel Uncertainty-Aware EWS under Discontinuous Telemetry, integrating a Split Conformal Prediction calibration layer over an $n-3$ historical spatial window and $L_2$-regularized Logistic Regression classifier. This conformal layer explicitly incorporates a missed-polling-cycle covariate to dynamically widen prediction sets as telemetry degrades, guaranteeing marginal coverage ($\ge 1-\alpha = 0.90$) without relying on arbitrary imputation heuristics. Validated against a dataset of 205,584 observations from the San Francisco Municipal Transportation Agency (SF Muni), the framework achieves a high Precision of 0.95 and an F1-Score of 0.8234. Furthermore, by employing a zero-ETL, in-process Hybrid Transactional/Analytical Processing (HTAP) architecture via DuckDB, the pipeline establishes a processing latency of \SI{0.00042}{\ms} per record. This performance definitively demonstrates that robust, mathematically calibrated EWS can operate on lightweight edge infrastructure without resorting to computationally exorbitant Big Data clusters, paving the way for resilient and generalizable predictive analytics in modern transit operations.
\end{abstract}

\begin{IEEEkeywords}
Public Transit, Early Warning Systems, GTFS-RT, Digital Twins, Machine Learning, Conformal Prediction, Uncertainty Quantification.
\end{IEEEkeywords}
